% !TEX root = ../main.tex
\documentclass[../main.tex]{subfiles}
\graphicspath{{\subfix{../images/}}}

\lstset{tabsize=2}

\begin{document}

\subsection{Реализация базовых механик}

\subsubsection{Движение игрока}

Функция \texttt{move\_player\_wasd} считывает состояние клавиш WASD или стрелок, формирует вектор направления и нормализует его перед применением скорости. Нормализация необходима для того, чтобы движение по диагонали не было быстрее движения по осям. После обновления позиции она ограничивается границами мира (\texttt{WORLD\_SIZE}).

\begin{lstlisting}[caption={move\_player\_wasd (упрощённо)}]
	void move_player_wasd(Player& player) {
		float dt = GetFrameTime();
		Vector2 move = {0, 0};
		if (IsKeyDown(KEY_W) || IsKeyDown(KEY_UP))   move.y -= 1;
		if (IsKeyDown(KEY_S) || IsKeyDown(KEY_DOWN)) move.y += 1;
		if (IsKeyDown(KEY_A) || IsKeyDown(KEY_LEFT)) move.x -= 1;
		if (IsKeyDown(KEY_D) || IsKeyDown(KEY_RIGHT)) move.x += 1;
		if (move.x != 0 || move.y != 0)
		move = Vector2Normalize(move);
		player.pos.x += move.x * player.speed * dt;
		player.pos.y += move.y * player.speed * dt;
	}
\end{lstlisting}

\subsubsection{Коллизии и неуязвимость}

Проверка столкновения снаряда (круг) с игроком (прямоугольник) реализована в функции \texttt{circle\_rect\_collision}. Она использует оптимизированный алгоритм: сначала проверяются расстояния по осям, затем, при необходимости, угол прямоугольника.

Функция \texttt{update\_projectiles} обрабатывает движение всех снарядов, проверяет столкновения и управляет неуязвимостью. Если игрок получает урон, его здоровье уменьшается, а \texttt{invincible\_until} устанавливается на \texttt{level\_time + 0.5f}. В течение этого времени повторные попадания не наносят урон.

\begin{lstlisting}[caption={Фрагмент update\_projectiles}]
	bool dead = false;
	if (hit_count > 0 && player.invincible_until < level_time) {
		dead = player.hit(1);
		player.invincible_until = level_time + 0.5f;
	}
	if (dead) game_state = GameState::GAME_OVER;
\end{lstlisting}

Для удаления снарядов, вышедших за пределы мира, используется идиома \texttt{std::remove\_if} с последующим \texttt{erase}. Это обеспечивает линейную сложность и безопасное удаление.

\subsection{Реализация системы событий и планировщика}

\subsubsection{Структура AttackEvent}

Событие атаки привязано к времени, прошедшему с начала уровня, и хранит функцию, которая будет вызвана в момент срабатывания. Благодаря \texttt{std::function} можно передавать как лямбды, так и именованные функции.

\subsubsection{Планировщик в update\_level}

Для событийных уровней (\texttt{BEGINNING}, \texttt{WAY}) используется статический вектор \texttt{events}, заполняемый при инициализации уровня. Сортировка событий по времени (функция \texttt{sort\_attack\_events}) необходима, чтобы выполнять атаки последовательно.

\begin{lstlisting}[caption={Цикл обработки событий}]
	while (next_event < events.size() && level_time >= events[next_event].time) {
		float dt = level_time - events[next_event].time;
		events[next_event].action(dt, projectiles, player);
		next_event++;
	}
\end{lstlisting}

Уровень считается завершённым, когда все события обработаны \texttt{(next\_event == events.size())} и вектор \texttt{projectiles} пуст.

\subsubsection{Вспомогательные функции для генерации атак}

Для сокращения дублирования кода созданы функции:
\begin{itemize}
	\item \texttt{add\_beam\_events} — серия пуль вдоль луча,
	\item \texttt{spawn\_bullet\_line} — одновременная линия,
	\item \texttt{spawn\_circular\_burst} — круговой разлёт,
	\item \texttt{add\_beam\_burst} — взрыв лучами,
	\item \texttt{add\_fan\_attack} — веер лучей,
	\item \texttt{spawn\_targeted\_bullet} — прицельный снаряд.
\end{itemize}

Эти функции используют только стандартные средства Raylib и C++ и не зависят от глобального состояния.

\subsection{Создание уровня «Beginning»}

Уровень «Beginning» (инициализация в \texttt{level\_beginning\_init}) предназначен для демонстрации основных механик и обучения игрока. Последовательность атак:

\begin{enumerate}
	\item Две плотные линии пуль, движущиеся сверху вниз (с разрывом в центре).
	\item Широкая линия с щелями.
	\item Бомба, падающая в центр, после взрыва — круговая волна.
	\item Веер из трёх лучей (центр, лево, право).
	\item Встречные горизонтальные лучи (слева направо и справа налево).
	\item Прицельная пуля, выпущенная из случайной точки вверху.
	\item Веер из пяти лучей.
	\item Вторая бомба с комбинированным взрывом (лучи + линия).
\end{enumerate}

Каждая атака добавляется в вектор \texttt{events} с указанием времени старта. Для бомб используется двухэтапное событие: сначала создание бомбы, затем через расчётное время — её взрыв с порождением новых событий. В момент взрыва бомба удаляется из списка снарядов поиском.

\subsection{Создание уровня «Way»}

Уровень «Way» (функция \texttt{level\_way\_init}) содержит более сложные паттерны:

\begin{enumerate}
	\item Два расходящихся луча из центра верхней границы.
	\item Широкая линия сверху.
	\item Две бомбы, падающие одновременно: одна даёт круговую волну, другая~— взрыв лучами.
	\item Веер из пяти лучей, горизонтальный луч и линия снизу.
	\item Прицельная пуля и встречная линия снизу.
	\item Два встречных веера (слева и справа).
	\item Каскадная бомба: после падения~— круговая волна и диагональные пули, которые при пересечении в центре порождают вторую круговую волну.
\end{enumerate}

% Особенность уровня — динамическое добавление новых событий внутри обработчика бомбы (функция \texttt{add\_beam\_burst}). После добавления вызывается \texttt{sort\_attack\_events}, чтобы сохранить временной порядок.

\subsection{Режим отладки и вспомогательные функции}

Режим отладки включается клавишей \texttt{F3} и управляется глобальной переменной \texttt{debug\_mode}. В этом режиме:
\begin{itemize}
	\item На экран выводятся FPS, координаты игрока, счётчик попаданий, текущее время уровня.
	\item Отображается хитбокс игрока (прямоугольник).
	\item В меню становятся видимыми тестовые уровни (\texttt{EMPTY}, \texttt{TEST}, \texttt{TEST\_HP}).
\end{itemize}

Для преобразования мировых координат в экранные с сохранением пропорций реализованы функции \texttt{world\_to\_screen} (для точки и для размера). Отрисовка спрайта игрока использует \texttt{DrawTexturePro} с автоматическим масштабированием, чтобы сохранить пропорции текстуры при изменении размера окна.

\subsection{Основные трудности и их решение}

В ходе разработки возникло несколько проблем, которые затем были решены. В таблице~\ref{tab:difficulties} перечислены основные трудности и способы их преодоления.

\begin{table}[H]
	\centering
	\caption{Трудности и их решения}
	\label{tab:difficulties}
	\begin{tabularx}{\textwidth}{|X|X|}
		\hline
		\textbf{Проблема} & \textbf{Решение} \\
		\hline
		Неправильная обработка коллизий & Замена функции расчёта коллизий. \\
		\hline
		Неправильная сортировка событий при динамическом добавлении & Вынесение сортировки в отдельную функцию, вызов после добавления новых событий \\
		\hline
		Завершение уровня при паузе отладки (level\_time скачком) & Ограничение dt сверху (0.5 сек) \\
		\hline
		Статическое состояние уровней не сбрасывалось при рестарте & Добавлен флаг reset\_level, обнуление static-переменных \\
		\hline
		Хитбокс игрока не соответствовал текстуре & Обрезана текстура \\
		\hline
	\end{tabularx}
\end{table}

Дополнительные сложности:
\begin{itemize}
	\item \textbf{Удаление бомб при взрыве}: в обработчике взрыва необходимо идентифицировать бомбу в векторе \texttt{projectiles}. Использована проверка по радиусу и приблизительному положению.
	\item \textbf{Нормализация движения}: без нормализации скорость по диагонали в $\sqrt{2}$ раз выше, что нарушало баланс. Исправляется вызовом \texttt{Vector2Normalize}.
	\item \textbf{Преобразование координат}: изначально игровая область неправильно центрировалась при изменении размера окна. Добавлено вычисление смещения \texttt{offset\_x}, \texttt{offset\_y}.
\end{itemize}

Эти исправления позволили получить стабильную версию игры.

\end{document}